\section{Discussion}
\label{sec:discussion}

\subsection{The Model Limits the Gates}

The design rests on the expectation that a system which checks its evidence and its
claims would turn confident errors into stated gaps. With the 7B model, only half of
this expectation was met: wrong answers became far less frequent, but easy questions
were held back as readily as hard ones. Because the graph, prompts and settings were
kept identical when the model was exchanged, the two explanations can be separated.
With the stronger model, G1 was convinced about five times as often, far fewer claims
were rejected by G2, and the held-back questions were the hard ones. The gates
therefore behave as intended when the model behind them can judge sufficiency and read
tables and charts reliably, and a 4-bit 7B model cannot yet do so consistently. A
practical arrangement follows from this result: a small local model can read and
write, while the two judgements are delegated to a stronger model.

\subsection{Measuring Grounding}

IoU against a hand-drawn box was fixed as the grounding measure before the evaluation,
but it was found to compare two drawing conventions as much as the precision of a
proof. A cell box spans its column, while an annotator boxes the number, or a whole
row for a comparison. The pointing game is insensitive to these conventions and is
therefore reported next to IoU. In future annotation, the unit should be fixed in
advance, for instance whole cells for tables.

\subsection{Threats to Validity}

The judge is the same model as the evaluated local systems; its agreement with a human
grader was substantial, but it was checked on 25 answers only. The answer boxes were
drawn by one annotator who, as a co-author, knew the aims of the work. The best-effort
rule was introduced after a first test run, and the pointing game after the IoU
results; both are disclosed, and the original measures are reported. The study covers
150 questions from 31 computer-science papers, retrieval within a known paper, one
quantised model run once at temperature 0, and a multi-hop subset of ten questions.
Reference attribution, used in place of a hand check of each citation, undercounts
claims that correctly cite supporting text. Finally, the benefit of proof boxes for
human verification was argued from their precision and was not measured in a user
study.
